% Part: many-valued-logic
% Chapter: sequent-calculus
% Section: rules-and-proofs

\documentclass[../../../include/open-logic-section]{subfiles}

\begin{document}

\olfileid{mvl}{seq}{rul}

\olsection{నియమాలు, \usetoken{P}{derivation}}

కింది చర్చలో $\Gamma, \Delta, \Pi, \Lambda$లు
\tetoken{వాక్యాల}{!!{sentence}s} పరిమిత క్రమాలను సూచిస్తాయని అనుకుందాం.

\begin{defn}[సీక్వెంట్]
\emph{$n$-వైపుల సీక్వెంట్} అంటే కింది రూపంలోని వ్యక్తీకరణ:
\[
\Gamma_1 \nSequent \dots \nSequent \Gamma_n
\]
ఇక్కడ ప్రతి $\Gamma_i$ భాష~$\Lang L$లోని
\tetoken{వాక్యాల}{!!{sentence}s} పరిమిత క్రమం;
అది శూన్యంగా కూడా ఉండవచ్చు.
\sourcecorrection{OLTEMVLSEQ-003}{ఆంగ్ల మూలం
ప్రతి స్థానానికి $\Gamma_1$ అని రాసింది.
$n$-వైపుల సీక్వెంట్‌లో ఏ స్థానాన్నైనా
సూచించేందుకు ఇక్కడ $\Gamma_i$గా సరిచేశాం.}
\end{defn}

\begin{defn}[ప్రారంభ సీక్వెంట్]
\emph{$n$-వైపుల ప్రారంభ సీక్వెంట్} అంటే,
భాషలోని ఏ \tetoken{వాక్యం}{!!{sentence}} $!A$కైనా
$!A \nSequent \dots \nSequent !A$
రూపంలో ఉండే $n$-వైపుల సీక్వెంట్.

భాషలో $0$-స్థానిక సంయోజకం~$\star$,
అంటే ప్రతిజ్ఞావాక్య స్థిరాంకం కూడా ఉంటే,
$\dots \nSequent \star \nSequent \dots$
అనే సీక్వెంట్‌నూ ప్రారంభ సీక్వెంట్‌గా
తీసుకుంటాం. ఇందులో $\star$,
$\tf{\star} \in V$కు సంబంధించిన
సత్యమూల్యపు స్థానంలో ఉంటుంది;
ఇతర స్థానాలు శూన్యంగా ఉంటాయి.
\end{defn}

$n$-విలువల తర్కం~$\Log{L}$లోని ప్రతి
సంయోజకానికి, $\Log{L}$లో ఆ సంయోజకం
పొందగల ప్రతి సత్యమూల్యానికి ఒక
తార్కిక నియమం ఉంటుంది. $\Log{L}$కు
చెందిన $n$-వైపుల సీక్వెంట్ కలనంలో
\tetoken{వ్యుత్పత్తులు}{!!^{derivation}s}
సీక్వెంట్ల వృక్షాలు. వాటి ఎగువన
ఉన్న సీక్వెంట్లు ప్రారంభ సీక్వెంట్లు;
ఒక సీక్వెంట్ మరొకటి లేదా మరికొన్నింటి
కింద ఉంటే, అది $\Log{L}$ సంయోజకాల
నిగమన నియమాన్ని సరిగ్గా వర్తింపజేసి
వాటి నుంచి రావాలి.

\begin{defn}[సిద్ధాంతాలు]
ఒక $n$-విలువల తర్కం~$\Log{L}$లో
\tetoken{వాక్యం}{!!^a{sentence}}~$!A$
\emph{సిద్ధాంతం} అవుతుంది,
$\Log{L}$లోని ప్రతి నిర్దేశిత
సత్యమూల్యానికి సంబంధించిన
స్థానంలో $!A$ ఉన్న $n$-వైపుల
సీక్వెంట్‌కు \tetoken{వ్యుత్పత్తి}{!!a{derivation}}
ఉంటే, అప్పుడు మాత్రమే.
$!A$ సిద్ధాంతమైతే
$\Proves[\Log{L}] !A$ అనీ, కాకపోతే
$\Proves/[\Log{L}] !A$ అనీ రాస్తాం.
\end{defn}

\begin{defn}[\tetoken{వ్యుత్పాద్యత}{!!^{derivability}}]
$n$-విలువల తర్కం~$\Log{L}$లో
\tetoken{వాక్యాల}{!!{sentence}s} సమితి~$\Gamma$
నుంచి \tetoken{వాక్యం}{!!^a{sentence}}~$!A$
\emph{\tetoken{వ్యుత్పాదించదగినది}{!!{derivable}}}
అని, $\Gamma \Proves[\Log{L}] !A$ అని అంటాం,
కింది షరతు నెరవేరితే, అప్పుడు మాత్రమే:
పరిమిత ఉపసమితి
$\Gamma_0 \subseteq \Gamma$ ఒకటి,
$\Gamma_0$లోని \tetoken{వాక్యాల}{!!{sentence}s}
క్రమం $\Gamma_0'$ ఒకటి ఉండి,
కింది సీక్వెంట్‌కు
\tetoken{వ్యుత్పత్తి}{!!a{derivation}}
ఉండాలి:
\[ \Lambda_1 \nSequent \dots \nSequent \Lambda_n \]
ఇక్కడ స్థానం $i$ ఒక నిర్దేశిత
సత్యమూల్యానికి సంబంధించినదైతే
$\Lambda_i$ అనేది $!A$;
లేకపోతే $\Gamma_0'$.
$!A$ను $\Gamma$ నుంచి
\tetoken{వ్యుత్పాదించలేకపోతే}{!!{derivable}}
$\Gamma \Proves/ !A$ అని రాస్తాం.
\end{defn}

ఉదాహరణకు, $3$-విలువల \L ukasiewicz
తర్కానికి $3$-వైపుల సీక్వెంట్ కలనం
ఉంది. $3$-వైపుల $\Gamma \nSequent \Pi \nSequent
\Delta$ అనే సీక్వెంట్‌లో $\Gamma$,
$\False$కు; $\Delta$, $\True$కు;
$\Pi$, $\Undef$కు సంబంధించినవి.
స్వీకృతాలు $!A \nSequent !A \nSequent !A$.
$\True$ ఒక్కటే నిర్దేశితమైనందున,
$\Gamma \Proves[\LogLuk[3]] !A$ అవుతుంది,
$\Gamma \nSequent \Gamma \nSequent !A$
అనే సీక్వెంట్‌కు
\tetoken{వ్యుత్పత్తి}{!!a{derivation}} ఉంటే,
అప్పుడు మాత్రమే. ($\Undef$ కూడా
నిర్దేశితమైతే, $\Gamma \nSequent !A
\nSequent !A$కు
\tetoken{వ్యుత్పత్తి}{!!a{derivation}}
అవసరమయ్యేది.)

\end{document}
